\section{Related Work} 
Network routing protocols are categorized into intra-domain and inter-domain routing. The application and scope of these protocols vary depending on network usage. This document examines the operation of Open Shortest Path First (OSPF) \citep{rfc2328}. OSPF is a link-state protocol that discovers network topology within a domain to facilitate packet transmission. However, OSPF has limited traffic control capabilities. Traffic engineering in OSPF is limited to tuning link cost weights. \citet{832225} discovered that OSPF link-weight adjustment cannot guarantee the optimal route because it is nondeterministic polynomial-time hard. Therefore, Equal-Cost Multi-Path (ECMP) is added to split traffic across equal-cost paths, load balancing traffic in the network. 

Researchers have studied different intra-domain protocols to improve network routing with traffic engineering.  \citet{article} compared protocols such as Multi Protocol Label Switching (MPLS), Open Shortest Path First (OSPF), and IP-based methods. Their findings show a trade-off between flexibility and scalability in traffic engineering. MPLS gives superior path control and supports traffic splitting. In contrast, IP-based and OSPF protocols are more scalable and cost-effective because these routing protocols do not require Label Switched Path (LSP) state and have an automatic failure re-routing mechanism.

To address these challenges, this study evaluates a recent protocol that uses machine learning as a baseline. In recent years, various machine learning approaches have been implemented in network routing. Reinforcement learning has become one of the leading approaches, where agents learn by trial and error during training to maximize system rewards \citep{SHAKARAMI2020107496, sutton_reinforcement_2018}. This approach even gained greater attention after the AlphaGo agent beat a human player \citep{shaheen2026reinforcementlearningstrategybasedatari}.

Several studies have explored reinforcement learning with neural networks to enhance routing performance. \citet{ALMASAN2022184} introduced a Deep Reinforcement Learning agent with Graph Neural Networks that generalizes across unseen topologies in optical networks.\citet{Yamin2024} applied Relational Actor-Critic to Integrated Access Backhaul networks, leveraging base-station relationships to group similar destinations. \citet{10225785} used a Multi-Agent Reinforcement Learning approach with Deep Q-Networks to reduce delays in Segment Routing networks. \citet{11134361} created a Multi-Agent Deep Reinforcement Learning method for dynamic routing in Mobile Ad Hoc Networks with Graph Neural Networks, providing a scalable and efficient solution for rapidly changing networks.

However, all of these works are evaluated on varied topologies and datasets. \citet{ALMASAN2022184} tested their DRL agent on 232 real-world network topologies from the Topology Zoo dataset to demonstrate generalized policy performance on different topologies.  \citet{11364348} assessed their cooperative MARINA framework on the GEANT2 network topology and found that the framework meets the quality of service (QoS) requirements more than traditional protocols like OSPF and ECMP. \citet{10293208} conducted simulation experiments with Abilene and GEANT with measured traffic data to show the reliability of G-Routing.

Despite strong results, these approaches share limitations relevant to this study. Most are evaluated using analytical models rather than packet-level simulation. Analytical models cannot capture queueing, loss, and delay outcomes. While packet-level simulation captures realistic network behavior, it requires more computing power \citep{1443314}. Some of the studies tried to tested with real packet-level demand on optical network \citep{ALMASAN2022184}, Software Defined Network \citep{11364348}, segment-routing \citep{10225785}, or mobile ad-hoc networks \citep{11134361}, but none of them utilize real measured network traffic.

This work addresses that gap by evaluating Multi-Agent Reinforcement Learning with Graph Neural Networks against the deployed baseline on intra-domain routing, using SNDLib topologies and real measured traffic \citep{SNDlib10, OrlowskiPioroTomaszewskiWessaely2010} and NS-3 \citep{Riley2010} for packet-level measurement.